\scanpage{51}
\mspage{13}
\uline{NB} \deleted{Si $\alpha:\unsure{\Delta}_{r'}\longrightarrow\Delta_r$ est surjectif, alors
$\alpha^*:\deleted{\unsure{X|I}}I_{\Delta_r}\longrightarrow I_{\Delta_{r'}}$ \unsure{est un}}
\inserted{On aurait intérêt à éliminer les}
``redondances'', en se bornant dans la
définition de
\[
X_{\Delta_r}\subset X\times I_{\Delta_r}
\]
aux drapeaux ``non dégénérés'' $\in I_{\Delta_r}$,
i.e. aux suites \uline{strictement} croissantes. Cela
oblige alors à se restreindre aux op.
$\frac12$ simpliciales relatives à \deleted{\unsure{des}} $\alpha:\Delta_{r'}\longrightarrow\Delta_r$ qui
sont \uline{strictement} croissantes. La définition
de $X_{\Delta_0}$ n'est pas changée, tandis que
$X_{\Delta_1}$ se trouve remplacé par $X_{\Delta_1}\setminus\text{diagonale}$
de $X_{\Delta_0}\times X_{\Delta_0}$.

\uline{Réconstitution de} $X$ à partir du syst.
$\frac12$ simplicial $X_{\Delta_*}$. Considérons d'abord
le diagramme
\diagram{\begin{tikzcd}[column sep=small,row sep=large]
&X_{\Delta_1}\arrow[dl,"a"']\arrow[dr,"b"]&\\
(19)\quad X_{\Delta_0}\arrow[dr,dashed]&&X_{\Delta_0}\arrow[dl,dashed]\\
&X&
\end{tikzcd}}
je voudrais qu'il soit cocartésien, i.e.
qu'on récupère $X$ comme somme amalgamée
\[
(20)\qquad X\simeq X_{\Delta_0}\amalg_{X_{\Delta_1}}X_{\Delta_0}.
\]
\mspage{50}
